\subsection{Source audit and compatibility at 236 states}

The newly supplied 236-state construction uses the same source DAG as the
237-state checkpoint, with one closed reassociation:
\[
 (233,(122,115))\quad\longmapsto\quad((233,122),115).
\]
The three chunks have 245, 160 and 132 instruction occurrences and no escaping
break, continue or halt.  Ordinary sequential associativity proves the source
equivalence for arbitrary natural register values.  Fresh builder and
independent-checker replays confirm exactly 537 changed paths, unchanged
ordered instructions and guard successors, and source score
$(224,6,48)\mapsto(224,6,49)$.

The reduction comes from the binary helper graph.  The reassociation creates
the edge $177/1\to212$, completing the predecessor-bit detour that permits
$223\mapsto182$.  Source node count is unchanged and path depth increases,
yet the final binary state count decreases.  This is a concrete reason to
optimize source association together with runtime helper availability.

Our additional checker reconstructs the source control graph by literal binary
increment/bit-stripping followed by zero-padded trie decoding.  Its equality
check is independent of the incoming ordered-interval construction.  The
parent's affine register invariants consequently transport unchanged by
instruction index.

The six-leaf NIP source rewrite lies at a disjoint prefix from the new rotation.
Applying the two rewrites in either order gives identical leaf lists: 6,449
occurrences, 224 live nodes, zero run six and depth 49.  The inherited NIP
scratch-register hypothesis remains necessary.  This source compatibility
does not validate the naive binary graft: on the literal 236 table, the
unchanged $\mathrm{CI}_8$ instruction at \texttt{01010} returns its legal zero
frame in 636 steps, whereas the five-row attempted edit returns the wrong tape
in 648 steps.  No 235-state machine is obtained from this composition.
